\chapter{Límites y continuidad}\label{ch:g12:limcont}

Este capítulo extiende la noción de límite de las sucesiones a las funciones
de una variable real e introduce la \omterm{def:g12:limcont:continuity}{continuidad}. El resultado central es el
teorema de los valores intermedios, que convierte una información
cualitativa (``$f$ es \omterm{def:g12:limcont:continuity}{continua} y cambia de signo'') en la existencia de
soluciones de ecuaciones.

\section{Límites de funciones}

En todo el capítulo, $f$ es una \omterm{def:g11:func:function}{función} definida en un \omterm{def:g10:numbers:interval}{intervalo} o en una
\omterm{def:g10:numbers:interunion}{unión} de \omterm{def:g10:numbers:interval}{intervalos} $D \subseteq \R$.

\begin{definition}[Límite en el infinito]\label{def:g12:limcont:liminf}
Sea $f$ definida en un \omterm{def:g10:numbers:interval}{intervalo} $\intco{a}{+\infty}$ y sea $\ell \in \R$.
Decimos que $f$ \emph{tiende a $\ell$ en
$+\infty$}\index{límite!de una función}, y escribimos
$\lim\limits_{x\to+\infty} f(x) = \ell$, si todo \omterm{def:g10:numbers:interval}{intervalo} abierto que
contiene a $\ell$ contiene todos los valores $f(x)$ para $x$
suficientemente grande: para todo $\varepsilon > 0$ existe $A$ tal que,
para todo $x \geq A$, $\abs{f(x) - \ell} \leq \varepsilon$.

Decimos que $f$ tiende a $+\infty$ en $+\infty$ si para todo $M \in \R$
existe $A$ tal que $f(x) \geq M$ para todo $x \geq A$. Los límites en
$-\infty$ y los límites iguales a $-\infty$ se definen de forma análoga.
\end{definition}

\begin{definition}[Límite en un punto]\label{def:g12:limcont:limpoint}
Sea $a \in \R$ y sea $f$ definida cerca de $a$ (salvo, quizá, en $a$).
Decimos que $f$ tiende a $\ell \in \R$ en $a$, y escribimos
$\lim\limits_{x \to a} f(x) = \ell$, si para todo $\varepsilon > 0$ existe
$\delta > 0$ tal que, para todo $x \in D$ con $\abs{x - a} \leq \delta$, se
cumple $\abs{f(x) - \ell} \leq \varepsilon$.

Decimos que $f$ tiende a $+\infty$ en $a$ si para todo $M$ existe
$\delta > 0$ tal que $f(x) \geq M$ siempre que $x \in D$ y
$\abs{x - a} \leq \delta$. Los límites laterales ($x \to a^+$,
$x \to a^-$) restringen $x$ a $x > a$ o a $x < a$.
\end{definition}

\begin{example}
$\lim\limits_{x\to+\infty} \frac{1}{x} = 0$,
$\lim\limits_{x\to 0^+} \frac{1}{x} = +\infty$,
$\lim\limits_{x\to 0^-} \frac{1}{x} = -\infty$. La \omterm{def:g11:func:function}{función}
$x \mapsto \frac1x$ no tiene límite (bilateral) en $0$.
\end{example}

\begin{definition}[Asíntotas]\label{def:g12:limcont:asymptote}
La recta $y = \ell$ es una \emph{asíntota horizontal}\index{asíntota} de la
curva de $f$ en $+\infty$ (resp.\ $-\infty$) si
$\lim_{x\to+\infty} f(x) = \ell$ (resp.\ en $-\infty$). La recta $x = a$ es
una \emph{asíntota vertical}\index{asíntota} si $f$ tiene un límite lateral
infinito en $a$.
\end{definition}

\begin{omfigure}
\begin{tikzpicture}
\begin{axis}[omaxis,
  width=8.8cm, height=6cm,
  xmin=-2.9, xmax=4.9, ymin=-5.2, ymax=9.2,
  xtick={1}, ytick={2}]
  \draw[omProp, dashed] (axis cs:-2.9,2) -- (axis cs:4.9,2);
  \draw[omProp, dashed] (axis cs:1,-5.2) -- (axis cs:1,9.2);
  \addplot[omDef, thick, domain=-2.8:0.855, samples=60] {2 + 1/(x-1)};
  \addplot[omDef, thick, domain=1.145:4.8, samples=60] {2 + 1/(x-1)};
\end{axis}
\end{tikzpicture}

{\small La curva de $f(x) = 2 + \frac{1}{x-1}$ con su \omterm{def:g12:limcont:asymptote}{asíntota horizontal}
$y = 2$ (en $\pm\infty$) y su \omterm{def:g12:limcont:asymptote}{asíntota vertical} $x = 1$.}
\end{omfigure}

\begin{proposition}[Operaciones con límites]\label{prop:g12:limcont:operations}
Las reglas de la \cref{prop:g12:seq:operations} (sumas, productos,
cocientes) valen literalmente para los límites de funciones, en un punto o
en el infinito, con las mismas cuatro indeterminaciones
$(+\infty)+(-\infty)$, $0\times\infty$, $\frac{\infty}{\infty}$ y
$\frac{0}{0}$.
\end{proposition}

\begin{proof}
Las demostraciones son idénticas al caso de las sucesiones, sustituyendo
``para $n \geq N$'' por ``para $x \geq A$'' o ``para
$\abs{x-a} \leq \delta$''.
\end{proof}

\begin{theorem}[Límite de una composición]\label{thm:g12:limcont:composition}
Sean $f, g$ funciones y $a, b, c$ \omterm{def:g10:numbers:sets}{números reales} o $\pm\infty$. Si
$\lim\limits_{x \to a} f(x) = b$ y $\lim\limits_{y \to b} g(y) = c$,
entonces
\[
\lim_{x \to a} g\bigl(f(x)\bigr) = c .
\]
El mismo enunciado vale con una \omterm{def:g12:seq:sequence}{sucesión} en lugar de $f$: si $u_n \to b$ y
$\lim_{y\to b} g(y) = c$, entonces $g(u_n) \to c$.
\end{theorem}

\begin{proof}
Tomemos $b, c$ finitos (los demás casos son análogos). Sea
$\varepsilon > 0$. Existe $\delta > 0$ tal que $\abs{y - b} \leq \delta$
implica $\abs{g(y) - c} \leq \varepsilon$; y existe $\delta' > 0$ tal que
$\abs{x - a} \leq \delta'$ implica $\abs{f(x) - b} \leq \delta$.
Encadenando las dos cosas se obtiene $\abs{g(f(x)) - c} \leq \varepsilon$
para $\abs{x - a} \leq \delta'$.
\end{proof}

\begin{theorem}[Comparación y sándwich]\label{thm:g12:limcont:squeeze}
Sean $f, g, h$ funciones y $a \in \R \cup \{\pm\infty\}$.
\begin{enumerate}
  \item Si $f \leq g$ cerca de $a$ y $f(x) \to +\infty$ cuando $x \to a$,
        entonces $g(x) \to +\infty$.
  \item Si $f \leq g \leq h$ cerca de $a$ y $f$ y $h$ tienden al mismo
        límite finito $\ell$ en $a$, entonces $g(x) \to \ell$ cuando
        $x \to a$.
\end{enumerate}
\end{theorem}

\begin{proof}
El mismo argumento que para las sucesiones
(\cref{thm:g12:seq:squeeze}).
\end{proof}

\begin{method}[Calcular límites en la práctica]\label{met:g12:limcont:limits}
\begin{itemize}
  \item Polinomios y funciones racionales en $\pm\infty$: saca como factor
        la mayor potencia de $x$; el límite es el del cociente de los
        términos de mayor grado.
  \item $\frac{0}{0}$ en un punto $a$: saca $(x-a)$ como factor del
        numerador y del denominador, o reconoce una \omterm{def:g11:deriv:derivative}{derivada}
        $\lim_{x\to a}\frac{f(x)-f(a)}{x-a} = f'(a)$
        (\cref{ch:g12:deriv}).
  \item Raíces cuadradas: multiplica por la expresión conjugada.
  \item Factores oscilantes ($\cos$, $\sin$): usa el teorema del sándwich.
\end{itemize}
\end{method}

\section{Continuidad}

\begin{definition}[Continuidad]\label{def:g12:limcont:continuity}
Sea $f$ definida en un \omterm{def:g10:numbers:interval}{intervalo} $I$ y sea $a \in I$. La \omterm{def:g11:func:function}{función} $f$ es
\emph{continua en $a$}\index{continuidad} si
$\lim\limits_{x \to a} f(x) = f(a)$. Es \emph{continua en
$I$}\index{continua} si es continua en todos los puntos de $I$.
\end{definition}

Intuitivamente, la \omterm{def:g10:functions:graph}{gráfica} de una \omterm{def:g11:func:function}{función} \omterm{def:g12:limcont:continuity}{continua} en un \omterm{def:g10:numbers:interval}{intervalo} se puede
dibujar sin levantar el lápiz.

\begin{proposition}[Continuidad de las funciones usuales]\label{prop:g12:limcont:usual}
Los polinomios, las funciones racionales, $x \mapsto \sqrt{x}$,
$x \mapsto \abs{x}$, $\cos$, $\sin$, $\exp$ y $\ln$ son \omterm{def:g12:limcont:continuity}{continuas} en todo
\omterm{def:g10:numbers:interval}{intervalo} contenido en su \omterm{def:g11:func:function}{dominio}. Las sumas, los productos, los cocientes
(donde estén definidos) y las composiciones de funciones \omterm{def:g12:limcont:continuity}{continuas} son
\omterm{def:g12:limcont:continuity}{continuas}.
\end{proposition}

\begin{proof}[Demostración parcial]
La estabilidad frente a sumas, productos, cocientes y composición se sigue
de los teoremas correspondientes sobre límites. La \omterm{def:g12:limcont:continuity}{continuidad} de
$x \mapsto x$ y de las funciones constantes es inmediata a partir de la
definición, así que los polinomios y las funciones racionales son
\omterm{def:g12:limcont:continuity}{continuos}. La \omterm{def:g12:limcont:continuity}{continuidad} de las demás funciones usuales se demuestra en
los capítulos donde se construyen, o bien se admite.
\end{proof}

\begin{example}
La \emph{\omterm{def:g11:func:function}{función} parte entera} $x \mapsto \floor{x}$ (el mayor entero
$\leq x$) no es \omterm{def:g12:limcont:continuity}{continua} en ningún entero $n$: su límite por la izquierda
allí es $n - 1$, mientras que $\floor{n} = n$.
\end{example}

\begin{omfigure}
\begin{tikzpicture}
\begin{axis}[omaxis,
  width=7.2cm, height=5.6cm,
  xmin=-1.9, xmax=3.5, ymin=-2.5, ymax=2.7,
  xtick={-1,1,2,3}, ytick={-2,-1,1,2}]
  \draw[omDef, thick] (axis cs:-1.5,-2) -- (axis cs:-1,-2);
  \draw[omDef, thick] (axis cs:-1,-1) -- (axis cs:0,-1);
  \draw[omDef, thick] (axis cs:0,0) -- (axis cs:1,0);
  \draw[omDef, thick] (axis cs:1,1) -- (axis cs:2,1);
  \draw[omDef, thick] (axis cs:2,2) -- (axis cs:3,2);
  \fill[omDef] (axis cs:-1,-1) circle (1.5pt);
  \fill[omDef] (axis cs:0,0) circle (1.5pt);
  \fill[omDef] (axis cs:1,1) circle (1.5pt);
  \fill[omDef] (axis cs:2,2) circle (1.5pt);
  \draw[omDef, fill=white] (axis cs:-1,-2) circle (1.5pt);
  \draw[omDef, fill=white] (axis cs:0,-1) circle (1.5pt);
  \draw[omDef, fill=white] (axis cs:1,0) circle (1.5pt);
  \draw[omDef, fill=white] (axis cs:2,1) circle (1.5pt);
  \draw[omDef, fill=white] (axis cs:3,2) circle (1.5pt);
\end{axis}
\end{tikzpicture}

{\small La \omterm{def:g11:func:function}{función} parte entera salta en cada entero: el valor (punto
relleno) es distinto del límite por la izquierda (punto hueco).}
\end{omfigure}

\begin{proposition}[Sucesiones y continuidad]\label{prop:g12:limcont:seqcont}
Si $u_n \to \ell$ y $f$ es \omterm{def:g12:limcont:continuity}{continua} en $\ell$, entonces
$f(u_n) \to f(\ell)$. En particular, si una \omterm{def:g12:seq:sequence}{sucesión} definida por
$u_{n+1} = f(u_n)$ \omterm{def:g12:seq:limit}{converge} a $\ell$ y $f$ es \omterm{def:g12:limcont:continuity}{continua} en $\ell$, entonces
$f(\ell) = \ell$.
\end{proposition}

\begin{proof}
La primera afirmación es el \cref{thm:g12:limcont:composition} aplicado con
la \omterm{def:g12:seq:sequence}{sucesión} $(u_n)$. Para la segunda, pasamos al límite en los dos miembros
de $u_{n+1} = f(u_n)$: el primero tiende a $\ell$ y el segundo, a
$f(\ell)$, y el límite es único.
\end{proof}

\section{El teorema de los valores intermedios}

\begin{theorem}[Teorema de los valores intermedios]\label{thm:g12:limcont:ivt}
\index{teorema de los valores intermedios}
Sea $f$ \omterm{def:g12:limcont:continuity}{continua} en $\intcc{a}{b}$ y sea $k$ un número cualquiera
comprendido entre $f(a)$ y $f(b)$. Entonces existe $c \in \intcc{a}{b}$ tal
que $f(c) = k$.
\end{theorem}

\begin{omfigure}
\begin{tikzpicture}
\begin{axis}[omaxis,
  width=8.8cm, height=6cm,
  xmin=-0.4, xmax=4.4, ymin=-0.5, ymax=4.7,
  xtick={0.5,3.19,3.8}, xticklabels={$a$,$c$,$b$},
  ytick={1.26,2.5,3.98}, yticklabels={$f(a)$,$k$,$f(b)$}]
  \addplot[omProp, dashed, domain=0:3.19] {2.5};
  \addplot[omDef, thick, domain=0.5:3.8] {0.08*x^3 - 0.5*x + 1.5};
  \draw[black, dashed, thin] (axis cs:0.5,0) -- (axis cs:0.5,1.26)
    -- (axis cs:0,1.26);
  \draw[black, dashed, thin] (axis cs:3.8,0) -- (axis cs:3.8,3.98)
    -- (axis cs:0,3.98);
  \draw[black, dashed, thin] (axis cs:3.19,2.5) -- (axis cs:3.19,0);
  \fill (axis cs:0.5,1.26) circle (1.5pt);
  \fill (axis cs:3.8,3.98) circle (1.5pt);
  \fill (axis cs:3.19,2.5) circle (1.5pt);
\end{axis}
\end{tikzpicture}

{\small El teorema de los valores intermedios: una curva \omterm{def:g12:limcont:continuity}{continua} que une
$(a, f(a))$ con $(b, f(b))$ tiene que cortar todas las rectas horizontales
$y = k$ situadas entre $f(a)$ y $f(b)$.}
\end{omfigure}

\begin{proof}
Sustituyendo $f$ por $f - k$ (y por $k - f$ si hace falta), podemos suponer
$f(a) \leq 0 \leq f(b)$ y buscar un cero de $f$. Construimos dos sucesiones
por \emph{dicotomía}\index{dicotomía}: tomamos $a_0 = a$, $b_0 = b$; dados
$a_n \leq b_n$ con $f(a_n) \leq 0 \leq f(b_n)$, sea
$m = \frac{a_n + b_n}{2}$ y definimos
\[
(a_{n+1}, b_{n+1}) =
\begin{cases}
(a_n, m) & \text{si } f(m) \geq 0,\\
(m, b_n) & \text{si } f(m) < 0.
\end{cases}
\]
En los dos casos, $f(a_{n+1}) \leq 0 \leq f(b_{n+1})$, la \omterm{def:g12:seq:sequence}{sucesión} $(a_n)$
es \omterm{def:g12:seq:monotonic}{creciente}, $(b_n)$ es \omterm{def:g10:functions:variations}{decreciente} y
$b_n - a_n = \frac{b-a}{2^n} \to 0$: las sucesiones son adyacentes (véase
el \cref{exo:g12:seq:10}) y convergen a un límite común
$c \in \intcc{a}{b}$. Por \omterm{def:g12:limcont:continuity}{continuidad}, $f(a_n) \to f(c)$ y
$f(b_n) \to f(c)$; pasando al límite en $f(a_n) \leq 0$ y en
$f(b_n) \geq 0$ se obtiene $f(c) \leq 0 \leq f(c)$, luego $f(c) = 0$.
\end{proof}

\begin{remark}
La demostración por \omterm{thm:g12:limcont:ivt}{dicotomía} es un algoritmo: da aproximaciones tan
precisas como se quiera de una solución, reduciendo el error a la mitad en
cada paso. Así es como un ordenador puede resolver $f(x) = 0$ sabiendo solo
evaluar $f$.
\end{remark}

\begin{theorem}[Teorema de la biyección]\label{thm:g12:limcont:bijection}
Sea $f$ \omterm{def:g12:limcont:continuity}{continua} y \emph{estrictamente \omterm{def:g12:seq:monotonic}{monótona}} en un \omterm{def:g10:numbers:interval}{intervalo} $I$.
Entonces, para todo $k$ estrictamente comprendido entre los límites de $f$
en los extremos de $I$, la \omterm{def:g10:algebra:equation}{ecuación} $f(x) = k$ tiene una \emph{única}
solución en $I$.
\end{theorem}

\begin{proof}
Existencia: elegimos $a, b \in I$ con $k$ entre $f(a)$ y $f(b)$ (es posible
porque los valores de $f$ se acercan a sus límites en los extremos) y
aplicamos el teorema de los valores intermedios en $\intcc{a}{b}$.
Unicidad: si $x < y$, la monotonía estricta da $f(x) \neq f(y)$, así que
dos puntos distintos no pueden ser los dos solución.
\end{proof}

\begin{method}[Resolver $f(x) = k$ con el teorema de la biyección]\label{met:g12:limcont:bijection}
Para demostrar que $f(x) = k$ tiene exactamente una solución en un
\omterm{def:g10:numbers:interval}{intervalo}:
\begin{enumerate}
  \item establece una \omterm{def:g10:functions:table}{tabla de variación} de $f$ (signo de $f'$, véase el
        \cref{ch:g12:deriv});
  \item en cada \omterm{def:g10:numbers:interval}{intervalo} donde $f$ sea \omterm{def:g12:limcont:continuity}{continua} y estrictamente \omterm{def:g12:seq:monotonic}{monótona},
        compara $k$ con los valores o los límites en los extremos;
  \item aplica el teorema de la biyección en cada uno de esos \omterm{def:g10:numbers:interval}{intervalos} y
        cuenta las soluciones.
\end{enumerate}
Una calculadora o la \omterm{thm:g12:limcont:ivt}{dicotomía} localizan después cada solución con la
finura que se quiera.
\end{method}

\begin{example}\label{ex:g12:limcont:cubic}
La \omterm{def:g10:algebra:equation}{ecuación} $x^3 + x = 1$ tiene una única solución real. En efecto,
$f(x) = x^3 + x$ es \omterm{def:g12:limcont:continuity}{continua} y estrictamente \omterm{def:g12:seq:monotonic}{creciente} en $\R$ (suma de
funciones estrictamente \omterm{def:g12:seq:monotonic}{crecientes}), $\lim_{-\infty} f = -\infty$ y
$\lim_{+\infty} f = +\infty$, así que el teorema de la biyección se aplica
con $k = 1$. Como $f(0.6) = 0.816 < 1$ y $f(0.7) = 1.043 > 1$, la solución
está en $\intoo{0.6}{0.7}$.
\end{example}

\section{Ejercicios}

\begin{exercise}[$\star$]\label{exo:g12:limcont:1}
Calcula los límites siguientes:
\[
\lim_{x\to+\infty} \frac{2x^3 - x + 1}{5x^3 + x^2}, \qquad
\lim_{x\to-\infty} \bigl(x^2 + 3x - 1\bigr), \qquad
\lim_{x\to 2^+} \frac{x+1}{x-2}, \qquad
\lim_{x\to+\infty} \frac{\cos x}{x}.
\]
\end{exercise}

\begin{exercise}[$\star$]\label{exo:g12:limcont:2}
Sea $f(x) = \dfrac{2x^2 - 3x + 1}{x - 1}$, definida para $x \neq 1$.
\begin{enumerate}
  \item Demuestra que $f(x) = 2x - 1$ para todo $x \neq 1$.
  \item Deduce $\lim\limits_{x \to 1} f(x)$. ¿Se puede prolongar $f$ a una
        \omterm{def:g11:func:function}{función} \omterm{def:g12:limcont:continuity}{continua} en $1$?
\end{enumerate}
\end{exercise}

\begin{exercise}[$\star$]\label{exo:g12:limcont:3}
Determina las \omterm{def:g12:limcont:asymptote}{asíntotas} de la curva de $f(x) = \dfrac{3x - 1}{x + 2}$ y
esboza la curva.
\end{exercise}

\begin{exercise}[$\star\star$]\label{exo:g12:limcont:4}
Calcula
\[
\lim_{x\to+\infty} \bigl(\sqrt{x^2 + x} - x\bigr)
\qquad\text{y}\qquad
\lim_{x\to 0} \frac{\sqrt{1+x} - 1}{x}.
\]
\end{exercise}

\begin{exercise}[$\star\star$]\label{exo:g12:limcont:5}
Sea $f(x) = x + \sin x$.
\begin{enumerate}
  \item Demuestra que $f(x) \to +\infty$ cuando $x \to +\infty$, aunque $f$
        no sea \omterm{def:g12:seq:monotonic}{monótona} en ningún \omterm{def:g10:numbers:interval}{intervalo} $\intco{A}{+\infty}$.
  \item Demuestra que $g(x) = \dfrac{x + \sin x}{x - \sin x}$ tiende a $1$
        en $+\infty$.
\end{enumerate}
\end{exercise}

\begin{exercise}[$\star\star$]\label{exo:g12:limcont:6}
Demuestra que la \omterm{def:g10:algebra:equation}{ecuación} $x^3 - 3x + 1 = 0$ tiene exactamente tres
soluciones reales y da, para cada una, un \omterm{def:g10:numbers:interval}{intervalo} de longitud $1$ que la
contenga. (Indicación: estudia la variación de $f(x) = x^3 - 3x + 1$; su
\omterm{def:g11:deriv:derivative}{derivada} es $3x^2 - 3$.)
\end{exercise}

\begin{exercise}[$\star\star$]\label{exo:g12:limcont:7}
Sea $f \colon \intcc{0}{1} \to \intcc{0}{1}$ \omterm{def:g12:limcont:continuity}{continua}. Demuestra que $f$
tiene un \emph{\omterm{pb:g10:functions:1}{punto fijo}}: existe $c \in \intcc{0}{1}$ con $f(c) = c$.
(Indicación: aplica el teorema de los valores intermedios a
$g(x) = f(x) - x$.)
\end{exercise}

\begin{exercise}[$\star\star\star$]\label{exo:g12:limcont:8}
Una senderista sube por un camino de montaña, sale a las 8:00 y llega a las
20:00. Al día siguiente baja por el mismo camino, saliendo otra vez a las
8:00 y llegando a las 20:00. Demuestra que hay una hora del día en la que
está exactamente en el mismo punto del camino los dos días. (Indicación:
introduce la \omterm{def:g11:seq:arithmetic}{diferencia} de las dos funciones de posición.)
\end{exercise}

\begin{exercise}[$\star\star\star$]\label{exo:g12:limcont:9}
Sea $f(x) = x^5 + x^3 + x - 1$.
\begin{enumerate}
  \item Demuestra que $f$ tiene un único cero real $\alpha$ y que
        $\alpha \in \intoo{0}{1}$.
  \item Describe un procedimiento de \omterm{thm:g12:limcont:ivt}{dicotomía} que calcule $\alpha$ con un
        error menor que $10^{-6}$ y da el número de iteraciones necesarias
        partiendo de $\intcc{0}{1}$.
\end{enumerate}
\end{exercise}

\section{Problema: No se cruza un río sin mojarse}

\begin{problem}[{Problema de fin de semana --- el teorema de los valores
intermedios impone puntos fijos, antípodas a la misma temperatura y mesas
que no cojean: tres teoremas de apariencia imposible a partir de una sola
idea}]
\label{pb:g12:limcont:1}
Arruga un mapa de un país y déjalo caer en cualquier punto de ese país:
algún punto del mapa queda \emph{exactamente} encima del lugar que
representa. En todo instante, dos puntos diametralmente opuestos del
ecuador están \emph{exactamente} a la misma temperatura. Y una mesa
cuadrada que cojea sobre un suelo irregular siempre se puede estabilizar
\emph{girándola}. Tres trucos de fiesta y un solo teorema: una magnitud
\omterm{def:g12:limcont:continuity}{continua} no puede pasar de negativa a positiva sin anularse
(\cref{thm:g12:limcont:ivt}). Este problema entrena los límites, construye
después las tres maravillas y sabe exactamente qué es lo que el teorema no
promete.

\medskip
\noindent\textbf{Parte I --- Soltura con los límites.}
\begin{enumerate}
  \item Calcula $\lim_{x \to +\infty} \dfrac{3x^2 - x}{x^2 + 5}$, \
        $\lim_{x \to +\infty} \dfrac{2x + 1}{x^2 + 1}$ \ y
        $\lim_{x \to +\infty} (x^3 - x^2)$
        (\cref{met:g12:limcont:limits}).
  \item Da las \omterm{def:g12:limcont:asymptote}{asíntotas} de $f(x) = \dfrac{3x - 1}{x + 2}$
        (\cref{def:g12:limcont:asymptote}), calculando los límites
        pertinentes.
  \item Calcula $\lim_{x \to 3} \dfrac{x^2 - 9}{x - 3}$ (factorizando) y
        $\lim_{x \to 0} \dfrac{\sqrt{x + 1} - 1}{x}$ (con el conjugado).
  \item Con el teorema del sándwich (\cref{thm:g12:limcont:squeeze}):
        calcula $\lim_{x \to 0} x \sin\frac1x$ y
        $\lim_{x \to +\infty} \frac{\sin x}{x}$.
  \item Por composición y términos dominantes:
        $\lim_{x \to +\infty} \dfrac{\sqrt{4x^2 + x}}{x}$.
\end{enumerate}

\medskip
\noindent\textbf{Parte II --- Raíces, certificadas.}
\begin{enumerate}[resume]
  \item Demuestra que la \omterm{def:g10:algebra:equation}{ecuación} $x^3 + x = 5$ tiene al menos una solución
        en $\intcc{1}{2}$.
  \item Demuestra que la solución es única en $\R$
        (\cref{thm:g12:limcont:bijection}).
  \item Demuestra el clásico: \emph{todo polinomio de grado \omterm{def:g11:func:parity}{impar} tiene al
        menos una \omterm{def:g11:quad:discriminant}{raíz} real}. (Escribe el argumento para
        $p(x) = x^3 + a x^2 + b x + c$: calcula los dos límites infinitos
        con el término dominante y aplica después el teorema de los valores
        intermedios en un \omterm{def:g10:numbers:interval}{intervalo} grande.)
  \item Localizar la \omterm{def:g11:quad:discriminant}{raíz} de la pregunta 6 por \omterm{thm:g12:limcont:ivt}{dicotomía}: ¿cuántas
        divisiones por la mitad de $\intcc{1}{2}$ garantizan un error menor
        que $10^{-6}$ (\cref{exo:g12:limcont:9})? Compara con la
        duplicación de cifras del método de Newton
        (\cref{pb:g11:deriv:1}): ¿cuántos pasos necesitaría más o menos?
  \item El contraejemplo que custodia el teorema: $g(x) = \frac1x$ cumple
        $g(-1) = -1 < 0$ y $g(1) = 1 > 0$ y, sin embargo, no se anula
        nunca. ¿Qué hipótesis del teorema falla, y qué enseña el ejemplo
        sobre comprobar las hipótesis?
\end{enumerate}

\medskip
\noindent\textbf{Parte III --- Tres teoremas de apariencia
imposible.}
\begin{enumerate}[resume]
  \item El teorema del \omterm{pb:g10:functions:1}{punto fijo} del segmento: si
        $f : \intcc{0}{1} \to \intcc{0}{1}$ es \omterm{def:g12:limcont:continuity}{continua}, demuestra que
        algún $c$ cumple $f(c) = c$. (Estudia $g(x) = f(x) - x$ en los dos
        extremos.)
  \item Traduce la pregunta 11 al mapa arrugado: ¿qué hace el papel de
        $f$, por qué aterriza en $\intcc{0}{1}$ (en sentido figurado:
        dentro del país) y cuál es el punto que queda exactamente sobre su
        propio lugar? (Se admite en silencio una sola dimensión: el
        enunciado honrado en dos dimensiones es el teorema de Brouwer, en
        los volúmenes universitarios.)
  \item El ecuador: sea $T$ una temperatura \omterm{def:g12:limcont:continuity}{continua} sobre una
        circunferencia, vista como una \omterm{def:g11:func:function}{función} $2\pi$-periódica con
        $T(0) = T(2\pi)$. Tomemos $g(t) = T(t) - T(t + \pi)$. Calcula
        $g(0) + g(\pi)$, deduce que $g$ se anula en algún punto de
        $\intcc{0}{\pi}$ y enuncia el teorema sobre puntos antípodas.
  \item La mesa que cojea: una mesa cuadrada perfectamente rígida se apoya
        sobre un suelo irregular pero \omterm{def:g12:limcont:continuity}{continuo} (¡sin escalones!); tres
        patas tocan siempre el suelo y una queda en el aire a una altura
        $h(\theta)$ que depende de forma \omterm{def:g12:limcont:continuity}{continua} del ángulo de giro
        $\theta$ de la mesa. Explica por qué girar un cuarto de vuelta
        obliga a la holgura a cambiar de signo y concluye que algún ángulo
        $\theta^*$ deja las cuatro patas apoyadas. (Este argumento se
        publicó como teorema de pleno derecho en 2005; los dueños de las
        cafeterías conocían el truco antes.)
  \item La senderista del \cref{exo:g12:limcont:8} usó el mismo esquema.
        Enuncia en tres líneas la plantilla que comparten las preguntas 11,
        13, 14 y la senderista: qué \omterm{def:g11:func:function}{función} construir, qué comprobar en los
        extremos y qué invocar.
  \item Lo que el teorema \emph{no} da: la existencia salió gratis, pero
        ¿qué dos preguntas deja abiertas y qué herramientas de este
        problema las responden (preguntas 7 y 9)?
\end{enumerate}

\medskip
\noindent\textbf{Parte IV --- Retratos con \omterm{def:g12:limcont:asymptote}{asíntotas}.}
\begin{enumerate}[resume]
  \item Reescribe $f(x) = \dfrac{3x - 1}{x + 2}$ como
        $3 - \dfrac{7}{x + 2}$ y úsalo para describir la curva por
        completo: \omterm{def:g12:limcont:asymptote}{asíntotas}, variación a cada lado y esbozo. (Es una
        \omterm{def:g11:func:reference}{hipérbola} desplazada; la curva del coste medio del
        \cref{pb:g10:reffunc:1} también lo era.)
  \item La \omterm{def:g11:func:function}{función} $f(x) = \dfrac{x^2 - 1}{x - 1}$ no está definida en $1$.
        ¿Qué valor asignado a $f(1)$ la hace \omterm{def:g12:limcont:continuity}{continua} allí, y por qué los
        analistas llaman a un punto así una discontinuidad
        \emph{evitable}?
  \item La \omterm{def:g11:func:function}{función} parte entera $x \mapsto \lfloor x \rfloor$ (el mayor
        entero $\leq x$): ¿dónde es \omterm{def:g12:limcont:continuity}{continua}, dónde salta y por qué falla
        para ella la conclusión del teorema de los valores intermedios
        (halla $a < b$ con
        $\lfloor a \rfloor < \frac12 < \lfloor b \rfloor$ y sin ninguna
        solución de $\lfloor x \rfloor = \frac12$)?
  \item Final: las dos caras de la \omterm{def:g12:limcont:continuity}{continuidad}: la promesa local (valor
        $=$ límite: nada de teletransporte) y el poder global (el teorema
        de los valores intermedios: se visitan todos los valores
        intermedios). Clasifica las cuatro maravillas de la Parte III bajo
        la cara que les corresponda y dale al eslogan del río su nombre
        matemático preciso.
\end{enumerate}
\end{problem}
